\begin{hebrewpage}[תקציר מנהלים]
\hebtitle{תקציר מנהלים}

\hebsubtitle{ClosetTwin: חדר הלבשה וירטואלי מבוסס AI}

ClosetTwin היא חברת B2B שמציעה לאתרי אופנה חדר הלבשה וירטואלי מבוסס AI: סריקת גוף מהטלפון, מידות הבגד ונטיית ההתאמה של הקונה מתחברות להמלצה על מידה. החזרות עולות לקמעונאים בארצות הברית בין 685 ל-890 מיליארד דולר ב-\range{2024}{2024}, לפי שתי הערכות שונות (Appriss ו-NRF). בסקר של Coresight הוחזרו \pct{24.4} מההזמנות המקוונות של ביגוד, ו-\pct{53} מהמשיבים ציינו מידה והתאמה כסיבה המרכזית. מחיר החזרה הוא כ-\pct{17} מעלות הייצור (במקרה של נעליים).

שוק חדרי ההלבשה הווירטואליים מוערך ב-8.5 מיליארד דולר ב-2025 (IMARC), אך זו הערכה של מו״ל ללא שיטה פומבית. המסחר המקוון בישראל מוערך ב-6.9 מיליארד דולר ב-2024, ואופנה היא הקטגוריה הגדולה בו. השוק הישראלי קטן: סך דמי השימוש האפשריים אצל קמעונאים מקומיים הוא כ-0.2 מיליון דולר בשנה, ולכן ישראל היא שוק הוכחה ואירופה, מהשנה הרביעית, נושאת את ההכנסות.

המחיר המוערך הוא 0.036 דולר להזמנה (כ-\ils{0.13}). הטענה המרכזית היא השערה, ולא עובדה: פיילוט אקראי של 4 שבועות יכול לגלות ירידה של כ-\pct{7.8} בהחזרות אצל קמעונאי ישראלי גדול. כלל הפרטיות העיקרי הוא תיקון 13 לחוק הגנת הפרטיות (בתוקף מאז 14 באוגוסט 2025): מזהה ביומטרי הוא מידע רגיש רק כשהוא משמש לזיהוי אדם, ו-ClosetTwin סורקת להתאמת מידה ולא לזיהוי. תקנות GDPR משניות ורלוונטיות בשלב אירופה.

אין רווח תפעולי בחמש השנים הראשונות. מסמך זה מתוארך ל-2026, והגיוס המבוקש הוא כ-2.8 מיליון דולר (ערך נגזר והנחה, שיתעדכן מהמודל). ערוץ המכירה הראשון הוא מכירה ישירה לקמעונאי אופנה מקומיים, ומשנה רביעית באירופה, גם דרך Shopify.
\end{hebrewpage}
